\subsection{Data Pipeline}
\label{sec:data-pipeline}

The data pipeline is an offline, batch-oriented process consisting of two stages. In the first, Python scripts and Jupyter notebooks clean and restructure the fourteen raw datasets: cell line names and gene symbols are mapped to their canonical DepMap and Ensembl identifiers (Section~\ref{sec::harmonisation}), DepMap's $\log_2(\text{TPM}+1)$ values are converted to linear units for comparability with HPA and GEO, gene-per-column matrices are pivoted into a row-per-observation long format keyed on \texttt{(ach\_id, ensembl\_id)}, and missing values are removed. The output is twelve harmonised CSV files: two dimension tables, a parent--child relationship table, three per-source expression fact tables, and six further fact tables covering proteomics, mutations, fusions, metabolomics, microRNA, and genomic signatures.

In the second stage, a conversion script reads each CSV and writes it to Apache Parquet format using PyArrow with Snappy compression. Frequently recurring string columns are dictionary-encoded as categoricals, boolean flags are cast to native types, and large tables are processed in 500{,}000-row chunks to limit peak memory. At application startup, the backend registers each Parquet file as a named view in an in-memory DuckDB instance, so that the data access layer can issue standard SQL without further file handling. DuckDB was chosen because it reads Parquet natively, requires no separate database server, and suits the single-node EC2 deployment. The pipeline is re-runnable: incorporating a new data release requires only re-executing the harmonisation scripts and conversion step, after which a service restart picks up the updated files. A metadata table records the release version and harmonisation date for reproducibility.


\subsection{Backend API}
\label{sec:backend-api}

The backend is a Python web service built with FastAPI. It exposes a stateless JSON API that the frontend calls over HTTP, organised into four router modules corresponding to distinct areas of functionality.

The \textbf{gene router} (\texttt{/api/genes}) provides a prefix-based autocomplete endpoint that the search bar queries as the user types a HUGO symbol, and a resolution endpoint that maps a symbol to its canonical Ensembl identifier. These two endpoints ensure that all downstream queries are keyed on stable Ensembl IDs rather than on display names that may be ambiguous.

The \textbf{ranking router} (\texttt{/api/rank}) accepts a \textsc{POST} request whose body specifies one or more target genes with an expression direction, RNA-to-protein weight, mutation and fusion filter modes, genomic instability thresholds, optional metabolite and miRNA exclusion criteria, disease and lineage filters, a tissue-boost flag, and the set of RNA sources to include. The endpoint resolves each gene symbol, delegates to the ranking service described in Section~\ref{sec::ranking-algorithm}, and returns the ranked list together with a copy of the query parameters so that the frontend can reproduce the request. Request and response payloads are validated through Pydantic schemas, which enforce type constraints, value ranges, and enumerated options (for example, restricting \texttt{mutation\_mode} to one of \texttt{ignore}, \texttt{include}, or \texttt{exclude}) before any computation begins.

The \textbf{cell line router} (\texttt{/api/celllines}) serves three endpoints. A detail endpoint returns the full evidence breakdown for a single cell line: per-source z-scores, percentile ranks, protein intensity, mutation and fusion annotations, and per-source combined scores computed by re-running the ranking pipeline with each RNA source in isolation. A similarity endpoint returns the most transcriptomically similar cell lines based on Pearson correlation across the DepMap expression profile. A comparison endpoint accepts a list of cell line identifiers and returns their expression, protein, and mutation data side by side for a nominated gene.

The \textbf{filter router} (\texttt{/api/filters}) returns the distinct values of disease, lineage, and subtype present in the cell line dimension table, along with the valid pairings between them, so that the frontend can populate cascading dropdown menus without embedding a copy of the vocabulary.

All four routers share a single data access layer that issues parameterised SQL against the DuckDB views described in Section~\ref{sec:data-pipeline}. Because Parquet categorical columns are stored as DuckDB \texttt{DICTIONARY} types, all string comparisons in the SQL layer explicitly cast to \texttt{VARCHAR} to avoid type-mismatch errors. FastAPI runs synchronous endpoint handlers in a thread pool; thread safety is maintained by creating a fresh DuckDB cursor for each query rather than sharing one across concurrent requests. The API is served by Uvicorn behind an Nginx reverse proxy on the production EC2 instance, with CORS configured to accept requests from the frontend origin.


\subsection{Frontend User Interface}
\label{sec:frontend-ui}

The frontend is a single-page application built with React and bundled by Vite. It communicates with the backend exclusively through the JSON API described in Section~\ref{sec:backend-api}, and all rendering state is held in React hooks so that the interface updates without full page reloads.

The layout is divided into a left sidebar for query construction and a right content area for results. In the sidebar, the \textbf{SearchBar} component provides a text input with debounced autocomplete: as the user types a HUGO gene symbol, the frontend queries the gene search endpoint and presents matching symbols with their Ensembl identifiers in a dropdown. Each gene is added together with a direction toggle (\textsc{high} or \textsc{low}), and the selected genes appear as colour-coded tags that can be individually removed. Below the search bar, the \textbf{FilterPanel} component exposes the full set of query parameters. Basic filters (disease, lineage, subtype) are presented as cascading dropdown menus whose option lists are fetched from the filter endpoints and narrowed automatically as the user makes selections: choosing a disease restricts the lineage dropdown to lineages recorded for that disease, and similarly for subtypes. When any tissue filter is active, a mode toggle allows the user to switch between hard filtering, which excludes non-matching cell lines entirely, and soft boosting, which retains all cell lines but promotes those matching the selected tissue. An expandable advanced section groups genomic filters (mutation and fusion include/exclude/ignore modes), scoring controls (RNA source toggles, an RNA-to-protein weight slider, a top-$N$ slider), genomic instability thresholds, metabolite and miRNA exclusion fields, and a planned assay selector. A badge on the toggle reports how many advanced filters are currently active.

Clicking ``Find Cell Lines'' sends the assembled parameters to the ranking endpoint. The \textbf{ResultsTable} component renders the response as a sortable table in which each row shows rank, cell line name, disease, lineage, confidence score (displayed as a coloured progress bar), evidence scenario, and source coverage count. Columns for mutations, fusions, assay compatibility, and tissue match are conditionally shown only when the corresponding filter is active, keeping the table uncluttered for simpler queries. If the user changes a filter after a search has been run, a visual indicator appears prompting them to refresh the results, and the refresh button is highlighted to signal the pending change.

Selecting a row opens the \textbf{DetailPanel}, a slide-in overlay that fetches the full evidence breakdown for that cell line: per-source z-scores, percentile ranks, per-source combined scores on the same scale as the overall score, protein data, and, when the relevant filters are active, mutation and fusion annotations. The panel also offers a ``Similar Cell Lines'' section that, on demand, retrieves the five most transcriptomically similar lines via the similarity endpoint, so the user can identify alternatives without running a new query. Users can select multiple rows and open a \textbf{CompareView} modal, which displays expression, protein, and mutation data side by side across the selected cell lines for each queried gene, with tabs to switch between genes in a multi-gene query.

The frontend is built to static assets via \texttt{npm run build} and served by Nginx on the same EC2 instance as the backend. In development, Vite's built-in proxy forwards \texttt{/api} requests to the local FastAPI server, so the frontend and backend can be developed independently on different ports.


\subsection{Deployment}
\label{sec:deployment}

The application is deployed on a single Amazon EC2 instance running Ubuntu. The backend runs as a \texttt{systemd} service under Uvicorn, which is restarted automatically if it crashes, and the frontend is compiled to static assets that Nginx serves directly. Nginx also acts as a reverse proxy, forwarding \texttt{/api} requests to the Uvicorn process on the local loopback interface. HTTPS is provided by a TLS certificate obtained through Certbot (Let's Encrypt) and renewed automatically, so that all traffic to \texttt{celllinefinder.com} is encrypted in transit.

Deployment is automated through a GitHub Actions workflow triggered on every push to the \texttt{mew/ranking-v2} branch. The workflow connects to the EC2 instance over SSH, pulls the latest code, installs any new Python or Node dependencies, rebuilds the frontend, and restarts the backend service. A health-check request to the \texttt{/health} endpoint runs after restart to verify that the backend is responding. This continuous-deployment pipeline ensures that a merged change reaches the live instance within minutes and that a failed deployment is surfaced immediately through the health check rather than discovered later by a user.